\documentclass{article}
\usepackage{graphicx} % Required for inserting images
\usepackage{algorithm}
\usepackage{algpseudocode}
\usepackage{amsmath} % 수학 기호 및 \text 사용을 위해 필요
\usepackage{amsmath,amssymb,amsthm}

\newtheorem{theorem}{Theorem}
\newtheorem{lemma}{Lemma}
\newtheorem{proposition}{Proposition}
\newtheorem{corollary}{Corollary}

\theoremstyle{definition}
\newtheorem{definition}{Definition}
\newtheorem{remark}{Remark}

\title{DDPC}
\author{yelimkim}
\date{December 2025}

\begin{document}

\maketitle

\section{Introduction}

\subsection*{Motivation}
Data-Enabled Predictive Control (DeePC) enables model-free predictive control directly from input--output data. 
However, in practical scenarios involving nonlinear or time-varying systems, 
online data selection becomes essential to ensure locality and computational tractability.

\subsection*{Need for Data Selection}
Data selection is typically motivated by:
\begin{itemize}
    \item \textbf{Locality:} selecting trajectory segments that are similar to the current operating regime.
    \item \textbf{Computation Time:} reducing the size of the Hankel matrix to lower online optimization cost.
\end{itemize}

\subsection*{Gap}
Existing selection-based DeePC approaches do not explicitly verify whether the selected data preserves 
\emph{quantitative persistency of excitation} (qPE), i.e., whether the minimum singular value condition remains satisfied after selection.

\subsection*{Contribution}
We propose a conditioning-aware data selection framework that:
\begin{itemize}
    \item Analyzes selection-induced degradation of quantitative PE.
    \item Connects conditioning loss to KKT sensitivity of the DeePC optimization.
    \item Introduces an adaptive-$K$ selection rule enforcing a minimum singular value margin.
\end{itemize}



\section{Preliminaries}

\subsection{Hankel Matrix Construction}
Define the Hankel matrices constructed from offline input--output trajectories.

\subsection{Willems' Fundamental Lemma}
State the classical result for LTI systems under persistency of excitation.

\subsection{Persistency of Excitation (PE)}
Discuss the rank-based PE condition.

\subsection{Quantitative PE}
Introduce the quantitative PE condition:
\[
\sigma_{\min}(H_u) \ge \gamma,
\]
which provides robustness in the presence of noise and uncertainty.

\subsection{Problem Statement}
Given a similarity-based selection mechanism, design a selection rule such that the selected Hankel matrix maintains quantitative PE.

\subsection{Notation}
Define symbols used throughout the paper.



\section{DeePC Formulation and KKT Sensitivity}

\subsection{Regularized DeePC as a Quadratic Program}
Present the regularized DeePC formulation with slack variables.

\subsection{KKT System}
Derive the associated KKT linear system.

\subsection{Sensitivity Bound}
Show that solution sensitivity is governed by the minimum singular value of the KKT matrix:
\[
\|\Delta x\|_2 \le \frac{1}{\sigma_{\min}(K)} \|\Delta b\|_2.
\]

\subsection{Relation to Singular Values of the Data Matrix}
Explain why conditioning degradation can be analyzed through singular values of the selected data blocks.



\section{Selection-Induced Conditioning Degradation}

\subsection{Similarity-Based Selection}
Define the baseline selection strategy.

\subsection{Singular Value Collapse}
Explain how similarity-based selection may induce redundancy and reduce $\sigma_{\min}$.

\subsection{Justification of the Proxy Matrix $A_g$}
Justify using
\[
A_g = \begin{bmatrix} U_p \\ Y_p \\ U_f \end{bmatrix}
\]
as a conditioning proxy instead of the full KKT matrix.



\section{Conditioning-Aware Adaptive Selection}

\subsection{Conditioning Margin}
Define the hybrid margin
\[
\Gamma = \max\{\gamma_{\min}, \rho \sigma_{\min}(A_{\rm ref})\}.
\]

\subsection{Adaptive-$K$ Rule}
Define
\[
K^\star = \min \{K : \sigma_{\min}(A_g(\mathcal{S}_K)) \ge \Gamma \}.
\]

\subsection{Termination Guarantee}
Provide a proposition ensuring termination under mild assumptions.

\subsection{Optional: Locality Metric Monotonicity}
Discuss locality--conditioning trade-off if needed.



\section{Numerical Examples}

\subsection{Evaluation Metrics}
Define phase transition metrics, slack-failure rate, success rate, and cost.

\subsection{Quantitative PE Collapse under Fixed $K$}
Demonstrate singular value collapse via $K$ sweep.

\subsection{Fixed-$K$ vs Adaptive-$K$}
Compare baseline selection and conditioning-aware selection 
in the inverted pendulum regulation task.

\subsection{Noise Sweep}
Evaluate robustness under increasing measurement noise.



\section{Conclusion}

Summarize the findings and outline directions for future research.

\end{document}